% SPDX-License-Identifier: LPPL-1.3c
% Copyright (C) 2026 Christophe Jorssen
% Author and Current Maintainer: Christophe Jorssen <christophe.jorssen@gmail.com>
% LPPL maintenance status: maintained
% This file is part of luafemm. See LICENSE and LICENSES.md for its terms.
\section{Rotors, stators and windings}
\label{sec:machines}

The machine nodes are genuine PGF shapes, using the same capture, materials,
mesh, solver and persistent-cache machinery as ordinary physical paths.
They describe a constant-depth planar cross-section with axial current density.
They require |field model=open|. Declaring a machine does not start a solve.
A rotor may precede or follow its stator: compound core contours preserve the
bore rather than painting an air disk over earlier regions.

\begin{keylist}[/tikz]{femm/rotor=\marg{options},femm/stator=\marg{options}}
Select a machine component and apply its options in |/femm/machine|.
The default argument is empty. Both keys enable |transform shape| and fix
node padding to zero. Physical dimensions are independent of node text.
Standard placement and transformation conventions are those of
Section~\ref{sec:components}. A circular gap scan additionally requires
concentric components with circular, unsheared transformed surfaces.
\end{keylist}

\subsection{Geometry and materials}
All lengths below are finite numerical literals in model millimetres; angles
are mechanical degrees, positive counterclockwise from local $+x$.
As in the existing component interface, dimension values are not unevaluated
PGF mathematical expressions. Evaluate expressions before supplying the keys.
Mesh sizes are measured in the model frame and do not scale with a node.

\begin{key}{/femm/machine/radius=\meta{positive number} (initially 29)}
Nominal rotor air-gap radius, before cutting its slots. Rotor only.
\end{key}
\begin{key}{/femm/machine/bore radius=\meta{positive number} (initially 30)}
Nominal stator bore radius, before cutting its slots. Stator only.
\end{key}
\begin{key}{/femm/machine/outer shape=\meta{circle or square} (initially circle)}
PGF choice selecting the stator core's outer outline. The bore remains circular.
\end{key}
\begin{keylist}[/femm/machine]{outer radius=\meta{number} (initially 45),
  outer width=\meta{number} (initially 90)}
Outer core radius for a circular stator, or full side length for a square
stator. The inactive dimension has no effect. The core must contain its bore,
slots and the requested minimum yoke.
\end{keylist}
\begin{key}{/femm/machine/shaft radius=\meta{nonnegative number} (initially 0)}
Rotor shaft radius. Zero omits the shaft region. A positive value cuts a
circular hole from the rotor core and fills it with |shaft material|.
\end{key}
\begin{key}{/femm/machine/housing thickness=\meta{nonnegative number} (initially 0)}
Separate stator casing outside the core, circular or square as selected above.
Zero omits it. It has its own material and drawing style.
\end{key}
\begin{key}{/femm/machine/minimum yoke=\meta{positive number} (initially 1)}
Minimum geometric clearance from slots to the shaft or outer core surface.
This is a construction check, not a saturation criterion. The angular-envelope
check also rejects overlapping adjacent slots conservatively; deeply staggered
slots with overlapping angular envelopes are not admitted.
\end{key}
\begin{keylist}[/femm/machine]{core material=\meta{material} (initially pure-iron),
 shaft material=\meta{material} (initially air),
 housing material=\meta{material} (initially air)}
Select catalogue identifiers or custom |femm/materials/| styles. The shaft and
housing have zero prescribed current. Conductors have independent selectors.
The default air housing is physically passive even when its style paints it.
\end{keylist}
\begin{key}{/femm/machine/magnetization angle=\meta{number} (initially 0)}
Local coercive-field direction for any coercive material used in a part.
It is transformed with the component. No automatic radial magnet segmentation
is implied by selecting a permanent-magnet material.
\end{key}
\begin{keylist}[/femm/machine]{mesh size=\meta{nonnegative number} (initially 0),
 slot mesh size=\meta{nonnegative number} (initially 0),
 conductor mesh size=\meta{nonnegative number} (initially 0)}
Local refinement for core/shaft/housing, slot cavities and conductors,
respectively. Zero inherits the problem spacing. Slot records can override
the latter two values. Refinement remains bounded by the ordinary mesh budget.
\end{keylist}
\begin{key}{/femm/machine/curve tolerance=\meta{positive number} (initially .03)}
Maximum chord sagitta used to discretise circular interfaces and notch fillets.
The arc step is also capped at ten degrees (fifteen for small fillets).
All machine material paths are emitted as polygons and are exported to FEMM
as segments. The nominal circles, their polygonal approximation and the mesh
are distinct: refine both tolerance and spacing in a convergence study.
\end{key}

\subsection{Slot placement and shape}
\begin{key}{/femm/machine/slots=\marg{options}}
Apply shared slot parameters in |/femm/machine/slots|. Its argument has no
implicit reset: PGF styles and grouping work normally. The defaults below
produce no slots until a count or explicit position list is provided.
\end{key}
\begin{key}{/femm/machine/slots/placement=\meta{uniform or explicit} (initially uniform)}
Choice between a generated angular sequence and the supplied |angles| list.
\end{key}
\begin{key}{/femm/machine/slots/count=\meta{integer} (initially 0)}
Number of slots for uniform placement, from 0 to 512. Zero makes a smooth
component. Explicit placement uses the number of entries in |angles| instead.
\end{key}
\begin{keylist}[/femm/machine/slots]{start angle=\meta{number} (initially 0),
 pitch=\meta{number} (initially 0)}
The $i$th uniform slot is at |start angle| plus $(i-1)$ times the pitch.
Zero pitch means $360/\texttt{count}$. Nonzero pitch allows a partial angular
array; angles are reduced modulo 360 for geometry. Coincident slots fail.
\end{keylist}
\begin{key}{/femm/machine/slots/angles=\marg{comma-separated numbers} (initially empty)}
Mechanical slot-centre angles for explicit placement, up to 512 entries.
Indices follow list order, even when the list is not sorted angularly.
\end{key}
\begin{key}{/femm/machine/slots/offset=\meta{number} (initially 0)}
Additional angular displacement, especially useful in an individual slot
record. It changes the geometry and the automatic winding's angular weight.
\end{key}
\begin{key}{/femm/machine/slots/shape=\meta{choice} (initially trapezoid)}
Choices are |rectangle|, |trapezoid| and |radial|. These describe the body
behind the narrow mouth. Rectangle has equal body widths; trapezoid uses the
independent bottom width; radial makes the body's two flanks pass through the
machine centre. The mouth itself has parallel sides.
\end{key}
\begin{keylist}[/femm/machine/slots]{depth=\meta{positive number} (initially 8),
 neck depth=\meta{positive number} (initially 1)}
Depth to the bottom, and depth to the shoulder where the mouth meets the
body. Both are measured along the slot's centre ray from the nominal surface,
into the iron: outward for the stator and inward for the rotor.
The total depth must exceed the neck depth.
\end{keylist}
\begin{keylist}[/femm/machine/slots]{opening width=\meta{positive number} (initially 2),
 opening angle=\meta{number} (initially 0)}
Width is the chord between the two mouth endpoints on the nominal circle;
it must be smaller than the nominal radius.
A positive opening angle overrides that width and specifies their subtended
angle; it must be less than 90 degrees. Zero selects the width convention.
This explicit precedence also applies inside styles and slot overrides.
\end{keylist}
\begin{keylist}[/femm/machine/slots]{width=\meta{positive number} (initially 4),
 bottom width=\meta{positive number} (initially 5)}
Tangential body width at the shoulder and bottom. The shoulder width must
contain the mouth. |bottom width| applies only to trapezoids; rectangles use
|width| and radial slots derive the bottom width from the two radii.
\end{keylist}
\begin{key}{/femm/machine/slots/corner radius=\meta{nonnegative number} (initially 0)}
Circular fillets on internal notch corners; zero leaves sharp corners.
The radius must fit adjacent edges and must not exceed conductor clearance.
Core and cavity share the same fillet vertices. The surface-mouth endpoints
remain on the nominal circle.
\end{key}
\begin{key}{/femm/machine/slots/material=\meta{material} (initially air)}
Material filling the entire cut slot before its conductors are inserted.
Use a custom nonmagnetic material to represent insulation. This region carries
zero imposed current. An insulating liner and residual air can be represented
by their common magnetic permeability in this magnetostatic model.
\end{key}
\begin{keylist}[/femm/machine/slots]{mesh size=\meta{nonnegative number} (initially 0),
 conductor mesh size=\meta{nonnegative number} (initially 0)}
Positive values override the machine's respective slot/conductor mesh keys.
Zero uses those machine values, which may themselves inherit the problem.
\end{keylist}
\begin{key}{/femm/machine/slot=\marg{options}}
Append an individual slot override. Supply |index|, a one-based slot number,
and any shared slot keys that describe its geometry, filling or conductors.
Unspecified fields retain the shared values. Placement, count, start angle, pitch and angles are global
array settings and are rejected in an individual slot record.
A slot index may appear only once. All records are compiled when the node is
constructed, after shared parameters are resolved, independently of the order
in which |slots| and |slot| appeared in that node's options.
\end{key}
\begin{key}{/femm/machine/slot/index=\meta{integer}}
Required existing slot index. Slot override fields live in |/femm/machine/slot|
and start empty; their names and meanings match the shared slot keys.
\end{key}

\subsection{Conductor sections and excitation}
\begin{key}{/femm/machine/slots/conductor material=\meta{material} (initially copper)}
Material of every conductor layer in the slot. This is independent of the
imposed current: changing copper to aluminium does not change a prescribed
$NI$. The present solver does not compute conductivity-driven eddy currents.
\end{key}
\begin{key}{/femm/machine/slots/clearance=\meta{positive number} (initially .4)}
Minimum inset from the straight body walls. Conductors occupy only the body,
not the mouth. The sloping side-wall inset is measured perpendicular to the
wall; the bottom and shoulder insets are radial. An excessive inset is an error.
\end{key}
\begin{keylist}[/femm/machine/slots]{layers=\meta{integer} (initially 1),
 layer gap=\meta{nonnegative number} (initially .3)}
One to eight equal-depth conductor layers, ordered from the shoulder towards
the bottom, with the stated radial separation. Each layer retains the local
trapezoidal width and has its own area. Unassigned layers are passive.
\end{keylist}
\begin{key}{/femm/machine/slots/current density=\meta{number} (initially 0)}
Uniform prescribed $J_z$ in A/m$^2$ in every conductor layer. Positive means
out of the page. A nonzero value conflicts with a coil or automatic winding
on the same layer. Raw density stays fixed under node scaling, unlike $NI$.
\end{key}
\begin{key}{/femm/machine/conductor=\marg{options}}
Append a conductor override, applied after slot geometry and before winding
assignment. Subkeys belong to |/femm/machine/conductor|: |slot=1|, |layer=1|,
|material|, |current density| and |style|. The last three start empty and leave
the corresponding slot setting unchanged. Repeated records apply in order.
Geometry is controlled by the slot's widths, depth, clearance and layer keys.
\end{key}
\begin{key}{/femm/machine/coil=\marg{options}}
Append a balanced two-sided coil. Subkeys are in |/femm/machine/coil|.
Multiple coils add their signed ampere-turns if they share a conductor layer.
The records do not draw end connections: the model is a transverse section.
\end{key}
\begin{key}{/femm/machine/coil/slots=\marg{first,second}}
Required two distinct existing slot indices. The first receives $+NI$ and
the second $-NI$. Their angular separation is unrestricted by the winding
interface; geometric collision checks still apply.
\end{key}
\begin{key}{/femm/machine/coil/layer=\meta{integer} (initially 1)}
Conductor layer used in both slots. It must exist in both.
\end{key}
\begin{keylist}[/femm/machine/coil]{turns=\meta{nonnegative number},
 current=\meta{number},ampere turns=\meta{number}}
Choose the product |turns*current| or the direct |ampere turns| value.
Defaults for the product are one turn and zero current. Specifying a direct
value together with either product key is rejected. Negative current or
ampere-turns reverses the coil. Fractional homogenised turns are permitted.
Actual polygon area determines $J_z=NI/S$ after physical transformations.
\end{keylist}
\begin{key}{/femm/machine/coil/material=\meta{material} (initially empty)}
Optional material override for both assigned conductor layers. Empty retains
their individual materials. If several coils override a shared layer's
material, the last such record wins.
\end{key}
\begin{key}{/femm/machine/winding=\marg{options}}
Configure an automatic distributed winding. Its subkeys are listed below.
Automatic distribution and explicit |coil| records cannot be combined.
\end{key}
\begin{key}{/femm/machine/winding/distribution=\meta{choice} (initially none)}
|none| leaves excitation to coil records or raw densities. |uniform| uses
weights $\operatorname{sign}\sin(p(\theta_i-\theta_0))$; |sine| uses
$\sin(p(\theta_i-\theta_0))$. Weights at numerical zeros are zero. Positions
must give a balanced sum. See Section~\ref{sec:tutorial-machines} for the
normalisation and its physical interpretation.
\end{key}
\begin{key}{/femm/machine/winding/ampere turns=\meta{number} (initially 1000)}
Total excitation assigned to the positive-weight slots, distributed across
all their layers. Negative values reverse the entire winding. This is not an
excitation per slot. The negative-weight slots carry the opposite total.
\end{key}
\begin{keylist}[/femm/machine/winding]{pole pairs=\meta{positive integer} (initially 1),
 axis=\meta{number} (initially 0)}
Spatial order $p$ and mechanical angular origin $\theta_0$ of the current
weights. Rotation of the whole node rotates its slots and sources together.
\end{keylist}

\subsection{Styles, anchors and a complete construction}
\begin{keylist}[/femm/machine]{core style=\marg{options},shaft style=\marg{options},
 housing style=\marg{options},slot style=\marg{options},conductor style=\marg{options}}
Append ordinary TikZ painting options to the respective part style. There are
no appended options initially. Physical changes belong in physical keys;
path replacements, transformations and physical-region keys in these painting
styles are unsupported.
\end{keylist}
\begin{keylist}[/femm/machine/slots]{style=\marg{options},conductor style=\marg{options}}
Extra painting options for a slot cavity or its conductor layers, empty
initially. Individual slot records can override them; |conductor| records
can override one layer's style. The global part style is applied first.
\end{keylist}
\begin{stylekey}{/tikz/femm machine core}
Initially |draw=black!70,fill=gray!15|.
\end{stylekey}
\begin{stylekey}{/tikz/femm machine conductor}
Initially |draw=orange!70!black,fill=orange!25|.
\end{stylekey}
\begin{stylekey}{/tikz/femm machine slot}
Initially empty; cavity material remains present even when invisible.
\end{stylekey}
\begin{stylekey}{/tikz/femm machine shaft}
Initially |draw=black!50|.
\end{stylekey}
\begin{stylekey}{/tikz/femm machine housing}
Initially |draw=black!70,fill=gray!30|.
\end{stylekey}

The standard compass anchors refer to the bounding box. |center|,
|core center| and |window center| coincide at the machine axis. The latter
also determines text placement. Physical anchors |gap east|, |gap north|,
|gap west| and |gap south| lie on the nominal air-gap circle. For each slot
$i$, the node supplies |slot i| (body centre), |slot i opening| (nominal surface),
|slot i bottom| and |slot i conductor j| (layer centre). Indices are literal
integers in the anchor name, for example |(S.slot 3 conductor 2)|.
An unavailable index raises an error when queried.

\luafemmexample{machine-nodes}

\subsection{Refinement and circular profiles between nodes}
\begin{key}{/tikz/femm/air gap=\marg{options}}
Set |/femm/air gap| options: required |rotor| and |stator| node names, and
|mesh size| (initially .5 mm, strictly positive). The |femm air gap| pic adds
an air refinement annulus after both nodes exist and before any solve.
It checks concentricity, circular transformed surfaces and positive clearance.
The annulus stays two curve tolerances from each nominal surface so that
polygonal approximation cannot paint air over iron. It refines the gap's
interior, not a prescribed number of radial element layers.
\end{key}
\begin{stylekey}{/tikz/femm air gap}
Initially empty. Paint the refinement band by redefining this style if desired.
\end{stylekey}
\begin{key}{/tikz/femm/gap profile=\marg{options}}
Set |/femm/gap profile| options for the |femm gap profile| pic:
|rotor| and |stator| are required names; |name=gap| names the registered
profile; |position=.5| is its fractional radial location from rotor to stator;
|samples=361| includes the closing endpoint; |angle origin=0| selects its
starting model-frame angle; |pole pairs=1| scales the electrical abscissa.
The position must lie strictly inside the gap, with clearance from the
polygonal walls. The pic registers an exact circular scan without solving.
Its visible annotation uses short straight segments. Its drawing transform is
the source model's frame; additional pic placement does not move the scan.
\end{key}
\begin{stylekey}{/tikz/femm gap profile}
Initially |draw=red!70!black,dashed,thin|. An empty drawing style can hide the
annotation without removing the profile.
\end{stylekey}

The gap must exceed four transformed curve tolerances. Uniform scaling and
rotation preserve circularity; unequal scaling and shear are rejected by the
gap helpers, though the individual nodes can still be transformed that way.
A gap-profile centre is retained with its source model, so subsequent figures
do not change its interpretation. For circulation, use an ordinary captured path rather than this exact circular
scan. Direct circular profiles are also available
as |require('luafemm-profiles').circle(model,name,x,y,radius,options)| in Lua.

\subsection{Polar plots and spatial harmonics}
\begin{keylist}[/femm/profile]{center x=\meta{number} (initially 0),
 center y=\meta{number} (initially 0),angle origin=\meta{number} (initially 0),
 pole pairs=\meta{positive integer} (initially 1)}
Polar reference for an ordinary captured profile. The centre is specified in
model coordinates. Components |Br| and |Btheta| use the outward radial and
counterclockwise angular unit vectors; |Hr| and |Htheta| follow the same rule.
At the centre, polar components are undefined and plotted as missing values.
Cartesian components retain their existing meaning.
\end{keylist}
The plot |component| choice accepts these four additional values. Its
|abscissa| choice additionally accepts |angle| and |electrical angle|.
Both are in degrees, relative to |angle origin|; the electrical angle is
|pole pairs| times the mechanical angle. Ordinary profiles unwrap successive
angles by the shortest increment, so sample densely enough to resolve the
winding around the centre. Exact circular profiles run from 0 to 360 degrees.

\begin{command}{\femmharmonicplot\oarg{options}}
Inside a PGFPlots axis, emit an ordinary |\addplot+| of harmonic amplitudes
against integer mechanical spatial order. No trailing semicolon is needed.
Requires a full circular gap profile (or Lua circular profile), with finite
samples. The repeated closing endpoint is omitted from Fourier sums.
\end{command}
\begin{command}{\femmharmonic\oarg{options}}
Print one spectrum statistic using PGF's number formatter. Both commands
parse their options locally in |/femm/harmonic| and can share family styles.
The calculation is pure Lua and reuses the profile's cached field samples.
\end{command}
\begin{key}{/femm/harmonic/profile=\meta{identifier}}
Required circular profile name; initially empty.
\end{key}
\begin{key}{/femm/harmonic/component=\meta{component} (initially Br)}
Field component to analyse: B, H, Az, their Cartesian, path-relative or polar
components. Circulation is not a periodic field component and is rejected.
\end{key}
\begin{key}{/femm/harmonic/maximum order=\meta{positive integer} (initially 15)}
Largest mechanical spatial order included. It must be strictly below half the
number of nonduplicate angular samples. This is a sampling limit, not a claim
that the FEM mesh resolves that harmonic. Increase both mesh resolution and
sampling and compare results when high-order amplitudes matter.
\end{key}
\begin{key}{/femm/harmonic/fundamental=\meta{positive integer} (initially 1)}
Reference order for normalisation and THD. For a $p$-pole-pair machine set it
to $p$. THD includes all other orders up to |maximum order|, including any
subharmonics, and excludes DC. It is a ratio, not a percentage.
\end{key}
\begin{key}{/femm/harmonic/normalize=\meta{true or false} (initially false)}
For spectrum plots, divide amplitudes by the selected fundamental. A zero
fundamental cannot be normalised. This key does not change scalar output.
\end{key}
\begin{key}{/femm/harmonic/order=\meta{positive integer} (initially 1)}
Order selected by scalar amplitude or phase output. It must not exceed
|maximum order|. Ignored for mean and THD.
\end{key}
\begin{key}{/femm/harmonic/quantity=\meta{choice} (initially amplitude)}
Scalar |amplitude|, |phase|, |mean| or |thd|. Phase is in degrees in the
convention $A\cos(n\theta+\phi)$, relative to the profile's angular origin.
A zero-amplitude phase is conventional, not physically meaningful.
Mean is the DC term; THD is undefined for a zero fundamental and raises an error.
\end{key}
\begin{key}{/femm/harmonic/plot options=\marg{PGFPlots options} (initially empty)}
Options passed to the emitted |\addplot+|. Axes, legends, colours, bar plots
and logarithmic scales remain standard PGFPlots settings.
\end{key}

Finite-element B is piecewise constant; plots deliberately retain those
jumps. Harmonic amplitudes are estimates from sampled FEM values, without
smoothing or forced sinusoidal fitting. Rotor-angle sweeps require distinct
static models and remeshing. Sliding air-gap elements, sector periodicity and
transient circuit coupling are separate extensions.
